\scanpage{7}\mspage{6}
Ainsi, l'ombre d'une figure est formée des lieux qui la raffinent -- on dit aussi que ce sont les \uline{lieux ``de la figure}'' (à ne pas confondre, avec les strates de la figure qui sont des lieux). La \unsure{\deleted{om}}multiombre est donc \unsure{formé|formée} des ombres de ses strates, elle est indexée par le déployement $\widetilde F$, ensemble des strates de $F$, la relation d'\deleted{\unsure{d'ordre|d'inclu}}\unsure{inclusion} de la figure ensembliste n'étant autre que la relation d'\unsure{incidences|incidence} entre strates de $F$, ie. la relation \unsure{de|dans} $\widetilde F$ induite par $\leq$ \unsure{de|dans} $\mathcal{M}$ (ou par $\ll$ au choix, ça revient au $\overline m$). Alors qu'une figure ne se réconstitue pratiquement jamais à partir de son ombre, elle se réconstitue pourtant \deleted{\unsure{à|d}}souvent (cas des ``\unsure{\deleted{a}}Ateliers ensemblistes'') \deleted{de}à partir de sa multiombre. Nous y reviendrons.

\S{} Il y a lieu de parler des relations d'ordre canoniques
